\chapter{The Completed Zeta Function and Its Fixed Axis}
\label{ch:zeta}

\section{Definition and analytic continuation}
For $\Ree s>1$, the Riemann zeta function is defined by the absolutely convergent Dirichlet series
\begin{equation}
\zeta(s)=\sum_{n=1}^{\infty}n^{-s}
\end{equation}
and by the Euler product
\begin{equation}
\zeta(s)=\prod_{p\ \mathrm{prime}}\left(1-p^{-s}\right)^{-1}.
\end{equation}
It extends meromorphically to the complex plane with a single simple pole at $s=1$. The nontrivial zeros lie in the critical strip $0<\Ree s<1$ and are symmetric about both the real axis and the critical line \cite{Riemann1859,Titchmarsh1986,DLMF25}.

A convenient entire completion is
\begin{equation}
\xi(s)=\frac12 s(s-1)\pi^{-s/2}\Gamma\!\left(\frac{s}{2}\right)\zeta(s).
\label{eq:xi-def}
\end{equation}
The factors outside $\zeta$ remove the pole and package the trivial zeros into the gamma factor. The completed function obeys
\begin{equation}
\xi(s)=\xi(1-s).
\label{eq:xi-fe}
\end{equation}
It also satisfies Schwarz reflection,
\begin{equation}
\xi(\overline{s})=\overline{\xi(s)},
\label{eq:xi-conj}
\end{equation}
because its defining data are real on the real axis and the analytic continuation respects conjugation.

\section{Linear and anti-linear symmetries}
There are two related transformations:
\begin{equation}
R(s)=1-s,
\qquad
J(s)=1-\overline{s}.
\end{equation}
Both are involutions:
\begin{equation}
R(R(s))=s,
\qquad
J(J(s))=s.
\end{equation}
The fixed points of $R$ are only the single complex point $s=1/2$, while the fixed points of $J$ form the entire vertical critical line.

\begin{theorem}[Fixed set of the zeta anti-involution]
For $s=\sigma+it$,
\begin{equation}
J(s)=s
\quad\Longleftrightarrow\quad
\sigma=\frac12.
\end{equation}
\end{theorem}

\begin{proof}
Since
\begin{equation}
J(\sigma+it)=1-(\sigma-it)=1-\sigma+it,
\end{equation}
the equation $J(s)=s$ is equivalent to $1-\sigma=\sigma$, hence $\sigma=1/2$. The imaginary coordinate remains arbitrary.
\end{proof}

Thus the critical line is not an arbitrary vertical line. It is the fixed locus of the reflection that combines the functional-equation complement $s\mapsto1-s$ with complex conjugation.

\section{Orbit structure of zeros}
If $\rho$ is a nontrivial zero, then the classical symmetries imply that
\begin{equation}
\rho,\quad \overline{\rho},\quad 1-\rho,\quad 1-\overline{\rho}
\end{equation}
are also zeros, with multiplicities preserved. If $\Ree\rho\neq1/2$ and $\Imm\rho\neq0$, these are generically four distinct points. If $\Ree\rho=1/2$, then $1-\overline{\rho}=\rho$, and the quartet collapses to the conjugate pair $\rho,\overline{\rho}$.

\begin{warning}
The functional equation guarantees orbit symmetry. It does not, by itself, require every zero to be a fixed point of $J$. Any proposed proof must supply a new reason why four-point orbits are impossible.
\end{warning}

This distinction is the logical center of the monograph. The information-spinor construction naturally selects fixed points, but a separate theorem is needed to show that all zero states belong to that fixed class.

\section{Centered coordinate}
Introduce
\begin{equation}
z=s-\frac12.
\end{equation}
Then the functional equation becomes an evenness relation. Define
\begin{equation}
\Xi_c(z)=\xi\!\left(\frac12+z\right).
\end{equation}
Equation \eqref{eq:xi-fe} gives
\begin{equation}
\Xi_c(z)=\Xi_c(-z).
\end{equation}
The critical line corresponds to purely imaginary $z=it$. The traditional real Riemann $\Xi$ function is
\begin{equation}
\Xi(t)=\xi\!\left(\frac12+it\right),
\end{equation}
which is real and even for real $t$.

The centered coordinate already has the form of a balance displacement. The real deviation
\begin{equation}
\delta=\sigma-\frac12
\end{equation}
changes sign under complement. The information model will identify this with the signed displacement from equal binary weight.

\section{The phase factor in the uncompleted functional equation}
One may write
\begin{equation}
\zeta(s)=\chi(s)\zeta(1-s),
\end{equation}
where one standard choice is
\begin{equation}
\chi(s)=2^s\pi^{s-1}\sin\!\left(\frac{\pi s}{2}\right)\Gamma(1-s).
\end{equation}
On the critical line $s=1/2+it$, the relation $1-s=\overline{s}$ and Schwarz reflection give
\begin{equation}
\zeta(s)=\chi(s)\overline{\zeta(s)}.
\end{equation}
Moreover $|\chi(1/2+it)|=1$. Therefore $\chi$ is a pure phase on the fixed axis. This established fact anticipates the later geometric statement: the critical line is where the dual branches are related by phase without a modulus imbalance.

Off the line, the duality factor generally changes modulus as well as phase. That observation motivates, but does not prove, the interpretation of $\sigma=1/2$ as the unitary or balanced duality axis.
